% AI-produced by GPT-6 Astra, OpenAI, 20 November 2026.
% Promoted as A.162 after complete root and independent proof/source audits.
% Derived from the preserved candidate; its historical SHA256 is recorded in A162_metadata.json.
\Needspace{10\baselineskip}
\subsection{A.162. An equivariant joint-germ model for projective concordances}

\begin{saresult}{A.162}{Proved under stated hypotheses}{The projective concordance fibre and a split forgetful kernel}

Let $D=D^4$ be the closed unit ball, $S=\partial D=S^3$, and let
$\iota(x)=-x$. Put
\[
 H=\operatorname{Homeo}_{\iota}(D),\qquad
 A=\operatorname{Homeo}_{\iota}(S),
\]
with their compact-open topologies. Every member of $H$ fixes the unique
fixed point $0$ of $\iota$. All groups below include every orientation
component.

Define a simplicial group $\mathcal G_{\iota}$ of \emph{joint germs} as follows.
A $q$-simplex is the germ along $\Delta^q\times\{0\}$ of a continuous
fibre-preserving homeomorphism between neighborhoods of that zero section
in $\Delta^q\times\mathbb R^4$, fixing the zero section pointwise and
commuting with $\iota$ in each fibre. Two representatives are equivalent
only when they agree on a neighborhood of the \emph{entire} zero section.
By compactness a representative and its inverse can be restricted to
uniform balls on the relevant side. The face and degeneracy maps are
parameter pullbacks. In particular, equality of two simplices is stronger
than pointwise equality of all their individual germs.

\textbf{Fibre model.}\quad
The restriction map
\[
 q:\operatorname{Sing}H\longrightarrow\mathcal G_{\iota}
 \tag{A.162.1}
\]
is a Kan fibration and is onto vertices. Its strict fibre over the identity
is weakly equivalent to $\operatorname{Sing}P_{\rm TOP}(\mathbb {RP}^3)$.
Boundary restriction and radial coning give $H\simeq A$. Consequently,
with the cone-germ map on the right, there is a natural weak equivalence
\[
 P_{\rm TOP}(\mathbb {RP}^3)
 \simeq_w\operatorname{hofib}\bigl(A\longrightarrow|\mathcal G_{\iota}|\bigr).
 \tag{A.162.2}
\]
Using the ordinary three-dimensional comparison and the generalized
Smale conjecture for $\mathbb {RP}^3$, this becomes
\[
 P_{\rm TOP}(\mathbb {RP}^3)
 \simeq_w\operatorname{hofib}\bigl(O(4)\longrightarrow|\mathcal G_{\iota}|\bigr)
 \simeq_w\Omega|\mathcal G_{\iota}/\operatorname{Sing}O(4)|.
 \tag{A.162.3}
\]
No topology on a set of individual germs is being substituted for the
specified simplicial model. In particular, no unproved comparison between
that model and a naive direct-limit topology is used. Let $\mathcal G_4$ denote the same joint-germ construction without the
oddness condition. The linear subgroup splits the forgetting map on
based fundamental groups, and there is a natural group isomorphism
\[
 \pi_0P_{\rm TOP}(\mathbb {RP}^3)
 \cong
 \ker\!\left(\pi_1|\mathcal G_\iota|
                 \longrightarrow\pi_1|\mathcal G_4|\right).
\]
In particular, the full component group is abelian. The kernel on the
right is not evaluated here. The statement does not identify the PL
image in it or prove zero-Casson--Sullivan sufficiency.
\end{saresult}

\begin{proof}
\textbf{Original inputs and their precise scopes.}\quad
The only equivariant existence input is Boyle--Chen--Conway,
\emph{Involutions on $S^4$}, arXiv:2512.22724v1 (27 December 2025),
Theorem~2.3, printed p.~6 \cite{BoyleChenConway2025}. It says that a free involution on $S^3\times I$
which does not swap its boundary components is conjugate to
$\iota|_{S^3}\times\operatorname{id}_I$. The proof uses Livesay's
classification of free involutions on $S^3$ and Kwasik--Schultz,
\emph{On $s$-cobordisms of metacyclic prism manifolds}, Inventiones~97
(1989), Theorem~2.1, printed p.~530 \cite{KwasikSchultz1989}, that the relevant topological
$\mathbb {RP}^3$ $h$-cobordism is a product. We use neither a PL product
conclusion nor a parametric extension from that theorem. The ordinary
four-dimensional annulus theorem is used in the same way as in their
Theorem~1.4, printed p.~15, case of zero-dimensional fixed-point set.

For parameter control we use Edwards--Kirby, \emph{Deformations of spaces
of imbeddings}, Annals of Mathematics~93 (1971), Theorem~5.1(1), printed
p.~70, and its explicitly canonical ambient-correction construction in the
proof of Corollary~1.2, printed p.~79 \cite{EdwardsKirby1971}. These apply to ordinary topological
manifolds, in particular the free quotient annulus. The deformation fixes
the inclusion and is supported in a compact neighborhood of the chosen
compact set. We derive the family version from this single continuous
operator and a finite subdivision of time; we do not cite a separate
parameterwise equivariant theorem.

\textbf{A normalized local extension operator.}\quad
Choose radii $0<a<b<c<d<1$ and the compact annular collar
$C=\{b\le\|x\|\le c\}$ in $D\setminus\{0\}$. Let $U$ be a slightly
larger compact annulus, contained in $\{a<\|x\|<d\}$, with $C$ in its
interior. Their quotients by $\iota$ are compact annuli in the manifold
\[
 M=(\operatorname{int}D\setminus\{0\})/\iota
       \cong\mathbb {RP}^3\times(0,1).
\]
Consider odd local homeomorphisms $f$ defined on a ball containing $U$,
fixing zero, and sufficiently close to the identity on $U$. We may require
that the image of that ball stay in $\operatorname{int}D$. On $\bar U=U/\iota$
the induced map $\bar f$ is an embedding into $M$ near the inclusion.

Edwards--Kirby gives a deformation $\varphi(\bar f,t)$ of that embedding,
constant at the inclusion, equal to $\bar f$ off a fixed compact
neighborhood of $\bar C$ in $\bar U$, and equal to the identity on
$\bar C$ at $t=1$. The construction used in their Corollary~1.2 produces
an ambient isotopy $\bar k_{f,t}$, starting at the identity and supported
in a fixed compact annulus after shrinking the embedding neighborhood,
whose terminal map agrees with $\bar f$ on $\bar C$.
Explicitly one uses
$\bar f\,\varphi(\bar f,t)^{-1}$ on the common image of $\bar U$ and the
identity outside. The images are the same since the deformation is fixed
near the boundary of the compact annulus. This is the same ambient
correction formula as in the cited proof.

Lift this isotopy through the double cover, starting from the identity.
It gives an odd ambient isotopy $k_{f,t}$ on $D\setminus\{0\}$, identical
near zero and the outer boundary, continuously depending on $f$. Extend
by the identity over these omitted regions. Its endpoint $k_f$ agrees
with $f$ on $C$: their lifts there differ by a constant deck transformation,
and the identity choice holds after shrinking the neighborhood, because
both maps at a fixed point of the connected collar are close to that
point. At $f=\operatorname{id}$ the full correction is the identity.

Choose a radius $r$ strictly between $b$ and $c$, and paste
\[
 E(f)(x)=
 \begin{cases}
 f(x),&\|x\|\le r,\\
 k_f(x),&\|x\|\ge r.
 \end{cases}
 \tag{A.162.4}
\]
The definitions agree on an annular neighborhood of the joining sphere.
They give a homeomorphism of $D$: both restrictions send the joining
sphere to the same locally flat sphere, and their inner images are the
component containing zero. The map is odd and agrees with $f$ near zero.
The compact-domain pasting gives joint evaluation continuity for families,
and inversion is continuous in the homeomorphism group of $D$.
Moreover $E(\operatorname{id})=\operatorname{id}$. Thus (A.162.4) is a
single normalized continuous extension operator on this neighborhood of
local maps. Its behavior inside the joining sphere is the given map,
not an independently selected filling.

\textbf{Every individual odd germ extends.}\quad
Given a germ $f$, choose $\epsilon>0$ so that $f$ is defined on a
neighborhood of $D_\epsilon$, with $f(D_\epsilon)\subset\operatorname{int}D$.
The complement
$W=D\setminus\operatorname{int}f(D_\epsilon)$ is topologically an annulus
by the four-dimensional annulus theorem. Its inherited involution is free,
and it preserves its two boundary components. Boyle--Chen--Conway's
Theorem~2.3 supplies an equivariant product parametrization of $W$.
Its parametrization on the inner boundary can be made exactly $f$ by
precomposing its $S^3$ coordinate with the required equivariant sphere
homeomorphism, held constant along the product coordinate. Glue this
parametrization to $f$ on $D_\epsilon$. The resulting odd homeomorphism
of $D$ extends the given germ; the outer boundary is allowed to move.
This proves vertex surjectivity of (A.162.1), with no family claim yet.

\textbf{Exact horn lifting with one common neighborhood.}\quad
Let $T=\Delta^j$, $j\ge1$, and let $Z$ be a simplicial horn. Suppose a
joint germ $g$ is given over $T$ and an actual continuous $H$-family
$F_z$, $z\in Z$, lifts its restriction. The finite horn data paste to a
continuous family on $Z$. Choose a retraction $r:T\to Z$ and a deformation
$r_t:T\to T$ with $r_0=r$, $r_1=\operatorname{id}$, and $r_t|_Z=\operatorname{id}$.
A horn is a strong deformation retract of a simplex, so such maps exist.

Extend the prescribed global family to $\widehat F_p=F_{r(p)}$.
Represent $g$ on a common ball and set, near zero,
\[
 \delta_p=\widehat F_p^{-1}g_p.
\]
Since equality on each horn face is equality of joint germs, finitely
many compact faces give a single smaller ball on which
$\delta_z=\operatorname{id}$ for every $z\in Z$. Restrict once to that
ball. The family
\[
 \delta_{p,t}=\delta_{r_t(p)}
 \tag{A.162.5}
\]
is an actual continuous local family, starts at the identity on one
common ball for every $p$, and stays the identity there for $p\in Z$.
The compact parameter set $T\times I$ and joint inverse continuity give
one still smaller ball contained in all the local images.

At times $t_i$, form the local increments
\[
 \eta_{p,t}=\delta_{p,t}\delta_{p,t_i}^{-1},
       \qquad t_i\le t\le t_{i+1}.
 \tag{A.162.6}
\]
They are defined on a fixed ball and equal the identity at $t_i$.
Use a fixed annular collar in that ball for the local operator (A.162.4),
rescaled spatially once. Compactness and continuity, including the
inverse evaluations, give a finite time partition such that every
increment (A.162.6) lies in the one extension neighborhood. This follows
by applying uniform continuity near the diagonal $t=t_i$ on the compact
set of all $(p,t_i,t)$; no new parameter-dependent choices are made.

Starting with $D_{p,0}=\operatorname{id}$, define successively
\[
 D_{p,t}=E(\eta_{p,t})D_{p,t_i}.
 \tag{A.162.7}
\]
At the joining times the factors $E(\operatorname{id})$ are identities,
so these families paste continuously. Their germs are exactly
$\delta_{p,t}$. On $Z$ all increments and all extension factors are
identities. Thus $\widehat F_pD_{p,1}$ is the required actual $H$-family,
with prescribed horn values exactly and germ $g$. This proves the Kan
fibration claim. The argument is a finite parameter argument and uses
joint germs throughout; an equality of isolated point germs would not
be sufficient for its common-ball step.

\textbf{The identity fibre and the concordance space.}\quad
Let $\mathcal F$ be the strict fibre of (A.162.1). Its simplices are actual
$D$-homeomorphism families equal to the identity on some common ball
$D_\epsilon$. Fix $\epsilon_0\in(0,1)$. The subgroup $H_{\epsilon_0}$
fixing $D_{\epsilon_0}$ pointwise is naturally homeomorphic to
$P_{\rm TOP}(\mathbb {RP}^3)$: quotient the outer annulus by $\iota$;
conversely lift a bottom-fixed projective concordance by its uniquely
specified identity lift on the inner boundary, and extend by identity
on the inner ball. Covering lifts and inverse evaluations are continuous
in every compact parameter family.

The inclusion $\operatorname{Sing}H_{\epsilon_0}\to\mathcal F$ is a
weak equivalence. A finite simplicial representative or homotopy in
$\mathcal F$ uses a common $\epsilon\le\epsilon_0$. Choose an odd radial
isotopy $a_t$ of $D$, fixing its boundary, whose radius map has graph
through
\[
 (0,0),\quad
 (\epsilon,(1-t)\epsilon+t\epsilon_0),\quad(1,1),
\]
affine on the two intervening intervals. It expands every radius and
sends $D_\epsilon$ to $D_{\epsilon_0}$ at $t=1$.
Conjugation by $a_t$ moves the given finite family into
$H_{\epsilon_0}$. It preserves $H_{\epsilon_0}$ throughout, because
$a_t(D_{\epsilon_0})\supset D_{\epsilon_0}$. This proves surjectivity on
homotopy groups. For injectivity, apply it to a null-homotopy with one
common $\epsilon$ and attach the boundary conjugation track, which lies
entirely in $H_{\epsilon_0}$. The same argument treats components.
Identity basepoints remain fixed. This identifies the full homotopy
fibre, without an incorrect contraction through merely pointwise germs.

\textbf{Replacing the middle group and identifying its map.}\quad
Restriction $H\to A$ has the continuous homomorphic section
\[
 c(a)(ru)=r\,a(u).
\]
For $h\in H$, the map $h c(h|_S)^{-1}$ fixes $S$ pointwise. Its kernel
contracts by the Alexander formula
\[
 \mathcal A_t(h)(x)=
 \begin{cases}t h(x/t),&\|x\|\le t,\ t>0,\\
 x,&\|x\|\ge t,
 \end{cases}
 \qquad \mathcal A_0(h)=\operatorname{id}.
\]
This is odd for every $t$, fixes zero and the outer boundary, and is
continuous in all compact parameters. Hence restriction is a homotopy
equivalence with inverse radial coning. Under this inverse, the map to
$\mathcal G_\iota$ is exactly the cone-germ map. Realizing the Kan fibre
sequence gives (A.162.2).

For completeness, the replacement of $A$ by $O(4)$ uses the
\emph{projective} Smale theorem, not the nonequivariant sphere theorem
alone. Bamler--Kleiner, \emph{Ricci flow and contractibility of spaces of
metrics}, arXiv:1909.08710, Theorem~1.3, printed p.~2
\cite{BamlerKleiner2019Metrics}, proves
$\operatorname{Isom}(M)\to\operatorname{Diff}(M)$ is a homotopy equivalence
for every round spherical quotient, including $\mathbb {RP}^3$.
Cerf, \emph{Groupes d'automorphismes et groupes de diff\'eomorphismes des
vari\'et\'es compactes de dimension 3}, BSMF~87 (1959), Theorem~5 and
Corollary~1, pp.~327--328 \cite{Cerf1959}, gives the ordinary Diff--TOP comparison;
its use of the sphere Smale assertion in Theorem~4, p.~327, is discharged
by the now-proved Smale theorem (also included in Bamler--Kleiner).
The full-proof account is Cerf, \emph{Topologie de certains espaces de
plongements}, BSMF~89 (1961), III, \S3.2.1, Theorem~8, p.~365, and
Corollary~1, p.~366 \cite{Cerf1961}; its sphere Smale hypothesis is
explicit and is discharged by the same input.

Every self-homeomorphism of $\mathbb {RP}^3$ lifts to $S^3$, and its two
lifts commute with the deck involution. The lift space is a double cover
of the homeomorphism group: nearby maps have a unique nearby lift after
specifying one sheet, by covering homotopy and compactness. The same holds
for diffeomorphisms and for the round isometries. Therefore the diagram
of double covers over
\[
 PO(4)\longrightarrow\operatorname{Diff}(\mathbb {RP}^3)
         \longrightarrow\operatorname{Homeo}(\mathbb {RP}^3)
\]
identifies its top-left group with $O(4)$ and its top-right group with
$A$. Both base maps are weak equivalences, hence their lifted maps are
weak equivalences as well. The resulting map to $\mathcal G_\iota$ is
the ordinary linear-germ inclusion. This proves the first equivalence
in (A.162.3). For the second, use the simplicial subgroup
$\operatorname{Sing}O(4)\subset\mathcal G_\iota$ and its Kan coset fibration.
The homotopy fibre of the fibre inclusion is the loop space of the base.

Finally $\pi_0O(4)\to\pi_0\mathcal G_\iota$ is a bijection. Surjectivity
follows from vertex surjectivity of (A.162.1) and $H\simeq A\simeq O(4)$;
orientation of an odd local homeomorphism separates the two components,
so injectivity holds. Consequently the exact sequence also gives
\[
 \pi_0P_{\rm TOP}(\mathbb {RP}^3)
 \cong
 \operatorname{coker}\bigl(\pi_1O(4)\longrightarrow
                              \pi_1|\mathcal G_\iota|\bigr),
 \qquad \pi_1O(4)\cong\mathbb Z/2.
 \tag{A.162.8}
\]
Here the group structure is the one induced by the fibre sequence of
simplicial groups, and it agrees with composition of projective
concordances under the fixed-radius fibre identification. The fundamental
group of a simplicial group is abelian. Hence (A.162.8) is an abelian group
cokernel, and in particular the component group on its left is abelian.

\textbf{Splitting by forgetting the odd symmetry.}\quad
Put $V=\pi_1|\mathcal G_\iota|$, $W=\pi_1|\mathcal G_4|$, and let
\[
 j:\pi_1O(4)\longrightarrow V,\qquad f:V\longrightarrow W
 \tag{A.162.9}
\]
be the linear inclusion and the forgetting homomorphism. The ordinary
joint-germ group here is literally the group $G$ defined in A.132,
with equality on a neighborhood of the whole parameter zero section.
Equation~(A.132.11) proves the weak equivalence between its identity
component and that of
$\operatorname{Sing}\operatorname{Homeo}(\mathbb R^4;0)$ by exact horn
lifting and the dilation contraction of the common-neighborhood fibre.
The map $e$ immediately following (A.132.12) agrees with the usual linear
action on $O(4)$. Thus this is a based comparison preserving the linear
map, rather than a substitution of individual germs or of an unrelated
germ-space topology.

The same A.132 calculation shows that
$\mathrm{TOP}(4)/O(4)$ has trivial $\pi_1$ and $\pi_2$: its stabilization
to $\mathrm{TOP}/O$ is $5$-connected and the latter two groups vanish.
The fibre exact sequence consequently makes
$\pi_1O(4)\to\pi_1\operatorname{Homeo}(\mathbb R^4;0)$ an isomorphism.
Under the preceding joint-germ bridge this map is exactly $fj$.
Therefore $j$ is split injective, and
\[
 V\cong j(\mathbb Z/2)\oplus\ker f,
 \qquad
 V/j(\mathbb Z/2)\xrightarrow{\ \cong\ }\ker f,
 \quad [v]\longmapsto v-j\bigl((fj)^{-1}f(v)\bigr).
 \tag{A.162.10}
\]
Combining this with (A.162.8) proves the asserted kernel formula.
\end{proof}

\textbf{Exact remaining comparison question.}\quad
Let $K=\ker f$ under (A.162.9), let
$I\subset K$ be the image of $\pi_0P_{\rm PL}(\mathbb {RP}^3)$ under the
proved identification, and let $\chi:K\to\mathbb F_2$ be the
Casson--Sullivan homomorphism $c_{\mathbb {RP}^3}$ of A.143 transported
through it. By A.142 this homomorphism is onto; by A.143 it vanishes on
$I$. Since $K$ is abelian, $I$ is a normal subgroup, and there is an
exact sequence
\[
 0\longrightarrow\ker\chi/I\longrightarrow K/I
       \xrightarrow{\ \chi\ }\mathbb F_2\longrightarrow0.
 \tag{A.162.11}
\]
Neither $K$, $I$, nor $\ker\chi/I$ is evaluated by A.162. In particular,
(A.162.11) does not assert $I=\ker\chi$, and it gives no assertion about
the kernel of the PL-to-TOP component map. This is a proved structural
reduction with no claimed movement of the original whole-space row or
of its zero-Casson--Sullivan sufficiency gate.

